\documentclass{osucourse}
\course{1125}

%% Not in the university's course data, so these come
%% from the published sheet this file was imported from.
\SetCourseField{exclusions}{Not open to students with credit for 106.}

\begin{document}

\CatalogDescription
\PrerequisitesAndExclusions

\begin{FollowupCourses}
Math 1126.
\end{FollowupCourses}

\begin{Purpose}
To develop an appreciation of, and basic competency in, the use of analytical thought in the development of a cohesive body of useful mathematical knowledge, with special emphasis on topics encountered in elementary and middle school mathematics programs. Math 1125 addresses the meaning of whole numbers, integers, rational numbers, and operations with these, number theory, and algebraic thinking. Appropriate only for those preparing to become early childhood educators and for those preparing to teach subjects other than math in middle school.
\end{Purpose}

\begin{Text}
Mathematics for Elementary Teachers, with Activity Manual, 4rd Edition, by Sybilla Beckmann, Pearson, ISBN for the package is 9780321836715 (loose-leaf) and Student Packet.
\end{Text}

\begin{TopicsList}
\item Counting and the decimal system.
\item Fractions and integers and their meaning.
\item Addition and subtraction of fractions, decimals, and integers.
\item Multiplication of fractions, decimals, and integers.
\item Division of fractions, decimals, and integers.
\item Ratios and proportional reasoning.
\item Number theory: factors and multiples, LCM, GCF, divisibility tests, prime numbers, unique factorization, notations for fractions and decimals.
\item Algebraic thinking: writing expressions, solving equations, sequences.
\item Problem solving and justification are themes of the course. *Currently taught in either lecture/recitation or workshop format.
\end{TopicsList}

\end{document}
